\section{The Hermitian--Einstein equation near the graded bundle}
\label{sec:construction}

We retain the hypotheses of Theorem~\ref{thm:main}. Norm estimates in
this section are uniform after shrinking the parameter neighborhood.
They do not compare $|s|$ with any positive slope surplus.

\subsection{Factor metrics and harmonic extensions}
Choose smoothly varying HE metrics $h_{i,s}$ on the fixed holomorphic
bundles $F_i$, with scalar normalizations. Here is the local analytic
justification. First choose normalized HE determinant metrics, using
the scalar Laplace equation. With those determinants fixed, the
linearization of the trace-free HE equation at $s=0$ is the
endomorphism Dolbeault Laplacian on trace-free Hermitian sections.
It is invertible because $F_i$ is stable. For an integer $\ell>n$,
the equation is a smooth map from $H^{\ell+2}$ to $H^\ell$.
The implicit-function theorem and elliptic regularity give the
metrics for small $s$. Their existence and simplicity also show that
the factors remain stable.

Let $\beta_{j,s}$ be the harmonic representative of the fixed nonzero
class in $H^{0,1}(\Hom(F_j,F_0))$. These representatives vary
smoothly. Indeed the holomorphic Hom bundle is fixed, so the
cohomology dimension is constant; the spectral projection onto the
kernel of its smoothly varying elliptic Laplacian is smooth. Elliptic
regularity gives smoothness in every spatial norm. Write
\[
 B_j(s)=\|\beta_{j,s}\|_s^2>0,\qquad
 h_{\mathrm{split},s}=\bigoplus_{i=0}^kh_{i,s}.
\]
Each $B_j$ has a uniform positive lower bound.

On the fixed smooth direct sum define
\begin{equation}\label{eq:scaled-dbar}
 \db_{s,t}=\db_0+\sum_{j=1}^kt_j\beta_{j,s},
 \qquad t\in\RR^k,
\end{equation}
where $\beta_{j,s}$ occupies the block $F_j\to F_0$.
The operator is integrable in every complex dimension:
$\db_0\beta_{j,s}=0$, and every product of two off-diagonal
entries vanishes. When all $t_j\ne0$, independent constant
scalings of the quotient factors identify this extension with
\eqref{eq:star}. Changes of harmonic representative are implemented
by upper-triangular complex gauges. Its determinant holomorphic
structure is independent of $t$.

\subsection{The complementary equation}
Let $\HH$ be the real space of constant block-scalar Hermitian
endomorphisms with total trace zero. It has dimension $k$. Let
$P_s$ be its $L^2_s$ projection and $Q_s=I-P_s$, acting on
pointwise trace-free Hermitian sections. We seek a metric
$h=h_{\mathrm{split},s}e^u$, where $u$ is a trace-free Hermitian
endomorphism satisfying $P_su=0$.
Express its mean curvature in the fixed Hermitian structure:
\begin{equation}\label{eq:Khat}
 \widehat K(s,t,u)
 =e^{u/2}\bigl(i\Lambda_{\omega_s}F_{\db_{s,t},h_{\mathrm{split},s}e^u}\bigr)e^{-u/2}.
\end{equation}
Because $\tr u=0$, its trace is the constant $r\kappa_{\omega_s}(E)$.

\begin{lemma}\label{lem:complement}
For small $(s,t)$ there is a unique small normalized $u(s,t)$ solving
\[
 Q_s\bigl(\widehat K(s,t,u)-\kappa_{\omega_s}(E)I\bigr)=0.
\]
It is smooth in the parameters and spatial variables, and for every $m$
\begin{equation}\label{eq:u-estimate}
 \|u(s,t)\|_{C^m}\le C_m|t|^2.
\end{equation}
The construction is equivariant under independent constant unitary
scalings of the factors.
\end{lemma}

\begin{proof}
Use the real Hermitian spaces $H^{\ell+2}\to H^\ell$, $\ell>n$.
Block-diagonal positive changes of metric identify the Hermitian bundles
for different $s$, preserve trace and block scalars, and are bounded
with bounded inverses in every fixed Sobolev norm. The projections
$Q_s$ depend smoothly on $s$. After those identifications,
$Q_s|_{\HH^{\perp_0}}$ is an isomorphism onto
$\HH^{\perp_s}$ for small $s$. These maps give fixed Banach spaces
for the parameter-dependent equation.

At $t=u=0$, the split curvature $K_{0,s}$ is block scalar, with
entries $\kappa_i(s)I_{F_i}$. Differentiation of the un-conjugated
mean curvature in the convention $h=h_{\mathrm{split},s}e^u$ gives
$\Delta_{\partial,s}u$. Differentiating the conjugation in
\eqref{eq:Khat} adds $[u,K_{0,s}]/2$. The K\"ahler identity on
endomorphism-valued zero-forms is
$\Delta_{\partial}-\Delta_{\db}=[K_{0,s},\cdot]$. Thus the real
Hermitian linearization is
\begin{equation}\label{eq:real-linearization}
 L_s=\Delta_{\db,s}+\tfrac12[K_{0,s},\cdot]
    =\tfrac12(\Delta_{\db,s}+\Delta_{\partial,s}).
\end{equation}
It is self-adjoint and nonnegative and annihilates $\HH$.
At $s=0$ all factor Einstein constants agree, so
$L_0=\Delta_{\db,0}$. Its trace-free Hermitian kernel is exactly
$\HH$: each factor has only scalar endomorphisms, and the Hom groups
between distinct stable factors of equal slope vanish.

Consequently $Q_0L_0$ is an isomorphism from the complementary
$H^{\ell+2}$ space to $H^\ell$. In fixed spaces, $Q_sL_s$ converges
in this operator norm to $Q_0L_0$. The Neumann-series criterion gives
uniformly bounded inverses. The implicit-function theorem now produces
$u(s,t)$, and uniqueness gives $u(s,0)=0$.

The derivative of the split mean curvature in $t_j$ has upper
off-diagonal entry $-\db_{\Hom}^*\beta_{j,s}$ and lower entry
its Hermitian adjoint. Both vanish by harmonicity for every $s$.
Differentiating the complementary equation therefore
gives $u_{t_j}(s,0)=0$, proving \eqref{eq:u-estimate}. Applying the
IFT at higher Sobolev indices and using uniqueness, or bootstrapping
the equation, gives each $C^m$ estimate. All equations, projections
and normalizations are invariant under the indicated unitary scalings;
uniqueness proves equivariance.
\end{proof}

\subsection{Scalar equations on the coordinate faces}
Set $p_j=I_{F_j}-(r_j/r)I$ for $1\le j\le k$. These form a basis
of $\HH$. The remaining curvature equations are
\[
 f_j(s,t)=\bigl\langle p_j,
 \widehat K(s,t,u(s,t))-\kappa_{\omega_s}(E)I\bigr\rangle_s=0.
\]
Changing the sign of any $t_i$ is induced by a unitary sign on $F_i$,
and fixes every $p_j$. Thus every $f_j$ is separately even in the
$t_i$.

We use the following elementary smoothness fact, jointly with the
parameters. A smooth function $f(s,t)$ that is separately even in
$t_1,\ldots,t_k$ can be written $F(s,t_1^2,\ldots,t_k^2)$ with
$F$ smooth on the closed positive orthant and smoothly extendible
across its faces. In one variable, Taylor expansion has only even
powers, and its differentiated remainders give derivatives of every
order in $w=t^2$ up to $w=0$. A smooth extension from the half-line
then gives the assertion; apply this successively in each variable,
retaining the other variables and $s$ as parameters. If such an $F$
vanishes on $w_j=0$, the formula
\[
 F(s,w)=w_j\int_0^1\partial_{w_j}F(s,w_1,\ldots,\theta w_j,
                                      \ldots,w_k)\,d\theta
\]
gives smooth divisibility on the orthant. We write $w_j=t_j^2$ below.

\begin{proposition}\label{prop:face}
The scalar equations have the form
\begin{equation}\label{eq:scalar-factor}
 f_j(s,w)=V_sr_j\epsilon_j(s)-w_jb_j(s,w),\qquad
 b_j(s,w)=B_j(s)+O\left(\sum_iw_i\right).
\end{equation}
The functions $b_j$ are smooth jointly in the parameters and orthant
variables, and the estimates are uniform near the origin.
\end{proposition}

\begin{proof}
On the face $t_j=0$, $F_j$ splits holomorphically from the rest.
Its circle of scalar unitary transformations fixes the Dolbeault
operator, the background metric and the complementary equation.
Lemma~\ref{lem:complement} implies that $u(s,t)$ commutes with
this circle. The corrected metric therefore also splits off $F_j$.
Conjugation by $e^{u/2}$ preserves the trace of that block. Since
the total residual is pointwise trace-free, Chern--Weil gives exactly
\begin{equation}\label{eq:face-identity}
 f_j(s,t)\big|_{t_j=0}
 =2\pi\deg_{\omega_s}(F_j)-V_sr_j\kappa_{\omega_s}(E)
 =V_sr_j\epsilon_j(s).
\end{equation}
This equality is independent of all other extension parameters.
The smoothness fact just stated gives the factorization in
\eqref{eq:scalar-factor}.

It remains to determine $b_j(s,0)$. At the split metric the quadratic
curvature term of the entry $t_j\beta_{j,s}$ is positive on $F_0$
and negative on $F_j$. Its pairing with $p_j$ is
$-t_j^2\|\beta_{j,s}\|^2$. The diagonal terms from other edges pair
to zero, as do their off-diagonal cross terms. The quadratic metric
correction contributes $L_su_2$, whose pairing with $p_j$ vanishes
because $L_s$ is self-adjoint and $L_sp_j=0$. Hence
$b_j(s,0)=B_j(s)$. Smoothness yields the remainder uniformly in $s$.
\end{proof}

\begin{proposition}\label{prop:balance}
For small $s$ with every $\epsilon_j(s)>0$, there is a positive
solution $w(s)$ of the scalar equations, and
\begin{equation}\label{eq:individual-balance}
 w_j(s)B_j(s)=V_sr_j\epsilon_j(s)\bigl(1+O(e(s))\bigr)
 \qquad(1\le j\le k).
\end{equation}
The estimates are uniform in the positive gap ratios. The corresponding
corrected metric is an HE metric on $E$, and $E$ is stable.
\end{proposition}

\begin{proof}
Extend the $b_j$ smoothly to a full neighborhood of $w=0$, preserving
positivity after shrinking it. Introduce independent real targets
$\eta_j$ and solve
\[
 w_jb_j(s,w)=V_sr_j\eta_j.
\]
The $w$-Jacobian at $w=\eta=0$ is $\diag(B_j(s))$, uniformly
invertible. The parameter IFT gives $w(s,\eta)=O(|\eta|)$,
with $w(s,0)=0$. For positive targets the equations themselves force
each $w_j>0$, so the solution lies where the extension agrees with
the original functions. Set $\eta_j=\epsilon_j(s)$. Then
\[
 \frac{w_jB_j(s)}{V_sr_j\epsilon_j(s)}
       =\frac{B_j(s)}{b_j(s,w)}=1+O(e(s)).
\]
No estimate divides an error by $\min_j\epsilon_j$.

Together with the complementary equation, vanishing of the scalar
pairings against a basis of $\HH$ gives the full HE equation. Take
$t_j=\sqrt{w_j}$; all are nonzero, so the holomorphic bundle is $E$.
To see simplicity, an endomorphism of $E$ preserves $F_0$ because
$\Hom(F_0,F_j)=0$. Its actions on $F_0$ and on each quotient
$F_j$ are scalar. Compatibility with each nonzero extension class
forces all these scalars to agree. The possible remaining map from
the quotient to $F_0$ is zero because $\Hom(F_j,F_0)=0$.
Thus an HE metric makes $E$ stable. Necessity of the positive
surpluses follows from the kernels of $E\to F_j$, as in the introduction.
\end{proof}
